\scanpage{344}
% top centre
\mspage{260}
% lines 1–16
Il reste seulement à prouver ce dernier complément $3^\circ$. C'est une contrainte \inserted{uniquement} sur la construction de $X_1(\cdot)$ et des $t_0(\lambda)$. Supposons donc que l'on ait construit $X_1(\cdot)$ sans tenir compte de cette contrainte. Pour $j$ fixé, il faut \unsure{savoir|donc} ``rendre commutatif\unsure{s|}'' tous ces carrés, i.e. toutes ces doubles-flèches $\bigl(t_0(\lambda)u(j'),\,\tau_0(j)u(j)F(\lambda)\bigr)$ de but $X_1(j)$. Comme l'ens de ces doubles-flèches est de cardinal $\leq\deleted{\unsure{\pi}}\unsure{\mathrm{card}|\mathrm{Card}}\operatorname{Ob}\unsure{\mathcal E|\mathscr E|J}_{/j}\leq\pi$, on peut trouver $X_1(j)\longrightarrow X'_1(j)$ dans $\mathcal E_j$, égalisant toutes ces doubles-flèches. \unsure{donc|dc} quitte à remplacer $X_1(j)$ par $X'_1(j)$, on \unsure{gagne|gagné}.
% archive sheet number: 343
